\section{COAST: a controlled cross-organ test}
\label{sec:protocol}
COAST follows three rules: the test organ cannot affect training or target selection; new patients
and new organs must be reported separately; and accuracy must be accompanied by information about
whether the target genes are measurable. These rules turn two public datasets into a reproducible
stress test.

Let a slide contain patch features $\mathbf Z\in\mathbb R^{N\times d}$, spot coordinates
$\mathbf C\in\mathbb R^{N\times2}$, and expression $\mathbf Y\in\mathbb R^{N\times g}$. A model
trained on organs $\mathcal O_{\rm tr}$ is evaluated on slides from either known organs or an organ
$o\notin\mathcal O_{\rm tr}$. All slides from a patient remain in one partition.

\subsection{Separate new patients from new organs}
\paragraph{LOOO: unseen-organ transfer.}
For each organ $o$, all cohorts and patients of $o$ are excluded from training and used only for the
final test. Model selection uses a fixed validation partition drawn from the remaining organs; the
test organ is not used for early stopping, hyperparameter selection, or threshold selection.

\paragraph{POOLED: known-organ deployment.}
All organs occur in training, but complete patients are held out within each organ. This evaluates a
single multi-organ model on new patients without conflating that question with unseen-organ transfer.

\begin{table}[t]
\centering\small
\setlength{\tabcolsep}{3.6pt}
\resizebox{\columnwidth}{!}{%
\begin{tabular}{lccc}
\toprule
Protocol & Patient held out & Organ held out & Primary question\\
\midrule
Within cohort & yes & no & in-distribution\\
\coast{}-POOLED & yes & no & multi-organ deployment\\
\coast{}-LOOO & yes & yes & organ transfer\\
\bottomrule
\end{tabular}
}
\caption{The protocols answer different questions and must be reported separately.}
\label{tab:regimes}
\end{table}

\subsection{Choose genes without looking at the test organ}
For fold $f$, we first form the technically compatible gene universe $A_f$ using assay annotations,
without inspecting held-out expression. We then rank genes by pooled log-transformed variance on
training slides only:
\begin{equation}
G_f=\operatorname{TopK}_{g\in A_f}
\operatorname{Var}_{(s,i)\in D^{f}_{\rm tr}}
\!\left[\log(1+Y_{sig})\right],\qquad K=50.
\label{eq:panel}
\end{equation}
The selected $G_f$ is frozen before validation and test evaluation. We release both $A_f$ and $G_f$
so that technical availability cannot be confused with expression-based selection. Clinically defined
genes are reported as a separate, predeclared analysis rather than inserted after inspecting results.

\subsection{Report when the target is measurable}
The primary endpoint is macro-averaged gene-wise Pearson correlation: calculate correlation over
spots for each nonconstant gene on each slide, average within slide, then give each organ equal weight.
We additionally report Spearman correlation, normalized RMSE, MAE, and detection average precision.
Micro (slide-weighted) means are secondary because HEST contains 2--24 slides per organ.

For every organ we report the fraction of spots with nonzero expression, median target variance, and
the number of evaluable genes. These are diagnostics, not inputs to the model. A detection-aware PCC
restricted to genes satisfying a threshold $\tau$ is a predeclared secondary endpoint; $\tau$ is fixed
without optimizing test performance and sensitivity is shown over several thresholds. Uncertainty is
estimated by a patient/slide bootstrap nested within organ and by three training seeds.

\subsection{Two complementary datasets}
HEST comprises 72 slides grouped into eight organs and spans Visium and Xenium; the distribution is
highly imbalanced. STImage-1K4M contributes 80 Visium slides, ten for each of eight organs, but its
released slides are heavily downsampled. Table~\ref{tab:data} reports the evaluation units.

\begin{table}[t]
\centering\small
\setlength{\tabcolsep}{3.5pt}
\resizebox{\columnwidth}{!}{%
\begin{tabular}{lrrr}
\toprule
Dataset / organ & Slides & Mean spots & Platform\\
\midrule
\multicolumn{4}{l}{\emph{HEST (72 slides)}}\\
Breast / Kidney / Colorectal / Prostate & 8/24/8/23 & 6938/3092/3007/2727 & mixed/Vis./Vis./Vis.\\
Lung / Pancreas / Liver / Skin & 2/3/2/2 & 2603/2524/2098/1517 & Xen./Xen./Vis./Xen.\\
\midrule
\multicolumn{4}{l}{\emph{STImage-1K4M (80 slides)}}\\
8 organs & 10 each & 942--3605 & Visium\\
\bottomrule
\end{tabular}
}
\caption{Cross-organ evaluation resources. Exact patient-level partitions are released as machine-readable files.}
\label{tab:data}
\end{table}

\subsection{Make every result traceable}
Each reported cell is tied to a versioned split file, panel file, prediction file, and evaluation
configuration. Code records patient identifiers after de-identification, training/validation/test
membership, seed, checkpoint selected on validation only, preprocessing version, and metric policy.
No test prediction is used to select a checkpoint or reporting threshold.
